% Generated draft: requires mathematical and source review.
\documentclass[11pt]{article}
\usepackage[margin=25mm]{geometry}
\usepackage{fontspec}
\setmainfont{Latin Modern Roman}
\usepackage{amsmath,amssymb,mathtools,unicode-math}
\setmathfont{Latin Modern Math}
\usepackage{enumitem,needspace,xurl,hyperref}
\setlength{\parindent}{0pt}
\setlength{\parskip}{5pt}
\setlength{\emergencystretch}{3em}
\newcommand{\N}{\mathbb{N}}
\newcommand{\Z}{\mathbb{Z}}
\newcommand{\Q}{\mathbb{Q}}
\newcommand{\R}{\mathbb{R}}
\newcommand{\C}{\mathbb{C}}
\newcommand{\F}{\mathbb{F}}
\newcommand{\A}{\mathbb{A}}
\newcommand{\PP}{\mathbb{P}}
\newcommand{\E}{\mathbb{E}}
\newcommand{\eps}{\varepsilon}
\newcommand{\dd}{\,\mathrm{d}}
\begin{document}
\section*{UnsolvedMath ID: 20002014}
\textbf{Source problem code:} catalogue entry
\subsection{Definitions and mathematical statement}
The source's actual Problem 1 asks for general or inductive formulae for operators whose extracted name ``§2`'' is corrupted. The aid recovers the name as the weighted tractor operator \(\square_{2\ell}\). The rest of the record begins a different problem's setup about affine conformal transformation laws and should not replace the operator-formula question.

Write \(\mathcal T[w]\) for a weighted standard tractor bundle and \(\mathcal E[w]\) for conformal densities of weight \(w\). In the usual weight convention for the family identified by the aid, its order-\(2\ell\) assignment is
\[
 \square_{2\ell}:\mathcal T[\ell-n/2]
              \longrightarrow\mathcal T[-\ell-n/2].
\]
At the Q-curvature order \(\ell=n/2-1\), where this order is defined, this specializes to
\[
 \square_{n-2}:\mathcal T[-1]\longrightarrow\mathcal T[1-n],
\]
whose bottom density slot is \(\mathcal E[-n]\), the slot of \(Q_n\) in the supplied recovery. The admissible dimensions, ambient extension, curvature corrections, normalization and exact operator definition are not supplied. The aid reports general strongly invariant tractor formulae, a curvature-corrected algorithm yielding a pure Q slot, low-order expansions and recursive Q-curvature formulae in primary literature, but explicitly says exact identification with the full ambient operator named in the AIM problem was not verified. Thus this is not a complete formula or a verified solution of Problem 1; the input producing Q and its corrections cannot be invented from the slot statement alone.

The accompanying affine-law discussion has the following precise independent content. On a fixed conformal class write \(\widehat g=e^{2\omega}g\), \(\omega\in C^\infty(M)\), and identify a geometric tensor-density target \(\mathcal V\) consistently across scales. Suppose for every representative and every \(\omega\) one has a finite law
\[
 N_{e^{2\omega}g}=N_g+L_g\omega,
 \qquad L_g:C^\infty(M)\to\Gamma(\mathcal V)\text{ linear}.
\]
Comparing two successive conformal changes gives
\[
 L_{e^{2\omega}g}\eta=L_g\eta
 \qquad(\omega,\eta\in C^\infty(M)).
\]
Thus \(L\) is conformally invariant in this density-target convention, and
\[
 C(g,e^{2\omega}g)=L_g\omega,\qquad
 C(g,g'')=C(g,g')+C(g',g'')
\]
defines an additive vector-valued cocycle. An infinitesimal first-variation law alone is not the assumed finite law.

If \(N_g\) is an integrable scalar density on a closed manifold and \(\int_MN_g\) is conformally invariant, then
\[
 \int_ML_g\omega=0\qquad\text{for every }\omega.
\]
Writing \(L_g\omega=(P_g\omega)\,\dd\mathrm{vol}_g\), this is \(P_g^*1=0\) for the formal adjoint. If additionally \(P_g\) is formally self-adjoint, it gives \(P_g1=0\). Self-adjointness is an extra assumption, not a consequence of the affine law. On a manifold with boundary or a noncompact manifold, support and boundary conditions are needed for this adjoint inference. For scalar Q-curvature, the finite density law is expressed with \(N_g=Q_g\,\dd\mathrm{vol}_g\); omitting the density identification from a scalar formula would lose its conformal weight factor.

Conversely, on one fixed conformal class choose a base metric \(g_0\), an initial target section \(N_0\), and a conformally invariant linear operator \(L\) into the fixed target. Then
\[
 N_{e^{2\omega}g_0}=N_0+L\omega
\]
satisfies the finite affine law. This construction uses a chosen base and does not by itself provide a natural local Riemannian tensor-density assignment independent of that choice. Locality and naturality are additional reconstruction requirements; conformal-path consistency is already ensured by the cocycle.
\subsection{Short English statement}
The damaged main question concerns tractor-operator formulae, whose exact ambient identification remains unverified. Its appended affine-law setup yields a separate compatibility theorem: the finite law forces operator invariance, while integral invariance and self-adjointness give different annihilation conditions. A base-dependent converse does not produce an intrinsic natural tensor.
\subsection{Sources}
\begin{itemize}
\item Permanent UnsolvedMath ID \(20002014\), exact public code catalogue entry; local registry \texttt{aim-geometry-notes.json}, \texttt{notes[351]}, canonical tag problem.
\item AIM workshop ``Conformal structure in geometry, analysis, and physics,'' Problem 1 followed by equation (34): \url{https://aimath.org/WWN/confstruct/confstruct.pdf}. No PDF was accessed. The aid's verification note reports using the hard-copy PDF and contemporaneous HTML ALT text to recover the corrupted operator, but the promised recovered sentence is absent from the supplied field.
\item Research attempt 1, ``Ambient tractor Laplacians and affine conformal laws,'' labelled partial\_result, partially\_solved, proved and low novelty confidence. It mentions primary tractor/Q-curvature literature without naming full references and expressly leaves exact full-ambient-operator identification unverified. Original literature, proposer/year metadata and full citations are absent.
\end{itemize}
\end{document}
